\section*{第 \ref{ch:g12:work-and-energy} 章——功与机械能}

\begin{solution}{exo:g12:work-and-energy:1}
$E_p = \tfrac12 \times 200 \times 0.050^2 = \qty{0.25}{J}$。拉伸量加倍，
能量为 $\times 4$，即 \qty{1.0}{J}。第二段 \qty{5.0}{cm}
单独需做功 $1.0 - 0.25 = \qty{0.75}{J}$——该段的力更大。
\end{solution}

\begin{solution}{exo:g12:work-and-energy:2}
$W = mg(z_A - z_B) = 0.90 \times 9.81 \times (12 - 30) \approx
\qty{-1.6e2}{J}$，与路径无关，沿直线爬升亦然。
往返一周的功为 \qty{0}{J}（保守力）。
\end{solution}

\begin{solution}{exo:g12:work-and-energy:3}
(a)、(b) 为保守力：功仅由端点决定（高度下降；
伸长量）。(c)、(e) 为非保守力：功为 $-fL$，阻力每行进一米
所做的功。(d) 不做功：力始终与运动方向垂直。
\end{solution}

\begin{solution}{exo:g12:work-and-energy:4}
$h = L(1 - \cos\qty{60}{\degree}) = \qty{1.0}{m}$;
$v = \sqrt{2 \times 9.81 \times 1.0} \approx \qty{4.4}{m/s}$。
绳的张力在每一时刻都与运动方向垂直：功率为零，
做功亦为零。
\end{solution}

\begin{solution}{exo:g12:work-and-energy:5}
$F = P/v = 250/9.0 \approx \qty{28}{N}$。速率为 \qty{5.0}{m/s} 时：
$F = 250/5.0 = \qty{50}{N}$（主要为斜坡重力的分量）。
\end{solution}

\begin{solution}{exo:g12:work-and-energy:6}
$E_p = \tfrac12 \times 450 \times 0.080^2 = \qty{1.44}{J}$;
$v = \sqrt{2 \times 1.44/0.020} = \qty{12}{m/s}$;
$h = E_p/mg = 1.44/(0.020 \times 9.81) \approx \qty{7.3}{m}$.
\end{solution}

\begin{solution}{exo:g12:work-and-energy:7}
梯形法（步长 \qty{2.0}{cm}）：$(0.03 + 0.10 + 0.19 + 0.30) =
\qty{0.62}{J}$。直线模型：$k = 18.0/0.080 =
\qty{225}{N/m}$，$\tfrac12 kx^2 = \qty{0.72}{J}$。实测曲线
在较小的 $x$ 处低于直线，因此真实面积
更小。
\end{solution}

\begin{solution}{exo:g12:work-and-energy:8}
$\Delta E_m = \tfrac12 \times 55 \times 9.0^2 - 55 \times 9.81 \times
8.0 = 2228 - 4316 \approx \qty{-2.1e3}{J}$;
$f = 2089/25 \approx \qty{84}{N}$。重力为保守力，法向力
不做功：只取决于下降高度和路径\emph{长度}，与
路径形状无关。
\end{solution}

\begin{solution}{exo:g12:work-and-energy:9}
$H_{\min} = \tfrac52 R = \qty{0.50}{m}$。从 \qty{60}{cm} 高处释放：
$v_{\text{最高点}} = \sqrt{2 \times 9.81 \times (0.60 - 0.40)} \approx
\qty{2.0}{m/s}$，且 $v_{\text{最高点}}^2/(gR) = 2.0$——为
最小值的两倍。摩擦会持续消耗能量，因此实际轨道需要更大的
高度。
\end{solution}

\begin{solution}{exo:g12:work-and-energy:10}
转折点：$8.0\,x^2 = 2.0$，$x = \pm\qty{0.50}{m}$。最大
速率在 $x = 0$ 处：$v = \sqrt{2 \times 2.0/0.100} \approx
\qty{6.3}{m/s}$。这是抛物线形势阱中的往复振动——
弹簧振子的 $E_p$ 恰好属于这种形式。
\end{solution}

\begin{solution}{exo:g12:work-and-energy:11}
$\tfrac12 kx^2 = \tfrac12 mv^2$:
$x = v\sqrt{m/k} = 2.0\sqrt{1200/\num{6.0e5}} \approx \qty{8.9}{cm}$.
速率为 \qty{4.0}{m/s} 时：\qty{18}{cm}——只是两倍而非四倍：能量
变为四倍，但 $E_p$ 随 $x^2$ 增长，因此 $x$ 只增为两倍。
\end{solution}

\begin{solution}{exo:g12:work-and-energy:12}
(a) $E_m = \qty{5.0}{J} < \qty{6.0}{J}$：它无法越过，将在
$E_p = \qty{5.0}{J}$ 处的斜坡停下并滚回。(b) 山顶：
$E_k = \qty{1.0}{J}$, $v = \sqrt{2 \times 1.0/0.50} = \qty{2.0}{m/s}$;
谷底：$E_k = \qty{5.0}{J}$，$v \approx \qty{4.5}{m/s}$。(c) $x = 0$ 处：
不稳定（极大值）；$x = \qty{1.5}{m}$ 处：稳定（极小值）。
\end{solution}

\begin{solution}{exo:g12:work-and-energy:13}
$\tfrac12 mv^2 = mgR_E$:
$v = \sqrt{2 \times 9.81 \times \num{6.37e6}} \approx \qty{1.12e4}{m/s}
= \qty{11.2}{km/s}$。各项都正比于 $m$：探测器与
小石子皆然。略小于此速度发射时，物体仍会上升极远，
到达转折点后落回，属于束缚运动。
\end{solution}

\begin{solution}{exo:g12:work-and-energy:14}
$E_m = \tfrac12 kA^2 = \tfrac12 \times 40 \times 0.060^2 =
\qty{7.2e-2}{J}$; $v_{\max} = \sqrt{2E_m/m} =
\sqrt{2 \times 0.072/0.250} \approx \qty{0.76}{m/s}$;
$E_k = E_p$ 的位置满足 $\tfrac12 kx^2 = E_m/2$：$x = \pm A/\sqrt 2 \approx
\pm\qty{4.2}{cm}$.
\end{solution}

\begin{solution}{exo:g12:work-and-energy:15}
$h = v^2/2g = 9.5^2/19.62 \approx \qty{4.6}{m}$；横杆高度 $\approx
4.6 + 1.0 = \qty{5.6}{m}$。纪录更高：运动员借助
手臂与腿在弯曲的撑竿上额外做功。但能量估算确定了
数量级——任何短跑选手的撑竿跳都达不到 \qty{10}{m}。
\end{solution}

\begin{solution}{pb:g12:work-and-energy:1}
\textbf{1.} $v = \sqrt{2gH} = \sqrt{2 \times 9.81 \times 45} \approx
\qty{29.7}{m/s}$：重力为保守力，法向力不做功——与
轨道形状无关。

\textbf{2.} 起滑点：$E_m = mgH = 70 \times 9.81 \times 45 \approx
\qty{3.09e4}{J}$。起跳点：$E_m = \tfrac12 \times 70 \times 26^2
\approx \qty{2.37e4}{J}$。

\textbf{3.} $\Delta E_m \approx \qty{-7.2e3}{J}$；
机械能定理把它归因于非保守力：
雪地摩擦力与空气阻力。

\textbf{4.} $f = \num{7.2e3}/90 \approx \qty{80}{N}$。

\textbf{5.} $\num{7.2e3}/\num{3.09e4} \approx 23\%$。涂蜡可减小 $f$，
蹲身可减小空气阻力：二者都使 $fL$ 减小，这是唯一可以调节的一项。

\textbf{6.} 水平初速度的自由落体——是一种
抛体运动。$v_x = \qty{26}{m/s}$；
$v_y = \sqrt{2 \times 9.81 \times 40} \approx \qty{28.0}{m/s}$。

\textbf{7.} $v = \sqrt{26^2 + 28.0^2} = \sqrt{1461} \approx
\qty{38.2}{m/s}$；能量法：$v = \sqrt{v_0^2 + 2gh} = \sqrt{676 + 785}$
——结果一致。约 \qty{138}{km/h}。

\textbf{8.} $t = v_y/g = 28.0/9.81 \approx \qty{2.9}{s}$；
$d = v_0 t \approx \qty{74}{m}$。

\textbf{9.} 不能。升力与速度垂直（不做功），
阻力与之相反（做负功）：$\Delta E_m \leq 0$，故实际落地速率只能
\emph{更小}——空气拉长了飞行过程，
而不是为其补充能量。

\textbf{10.} $\alpha = \arctan(28.0/26) \approx \qty{47}{\degree}$
速度方向低于水平方向。

\textbf{11.} 速度与斜面的夹角：$47 - 35 = \qty{12}{\degree}$。
$v_\parallel = 38.2\cos\qty{12}{\degree} \approx \qty{37.4}{m/s}$；
$v_\perp = 38.2\sin\qty{12}{\degree} \approx \qty{7.9}{m/s}$。

\textbf{12.} 被吸收的部分：$\tfrac12 \times 70 \times 7.9^2 \approx
\qty{2.2e3}{J}$，对比总动能 $\tfrac12 \times 70 \times 38.2^2 \approx
\qty{5.1e4}{J}$，两者之比约为 $4\%$。

\textbf{13.} $h_{\text{等效}} = 7.9^2/19.62 \approx \qty{3.2}{m}$——
相当于从墙上纵身一跳。平地上：$38.2^2/19.62 \approx \qty{74}{m}$——
无法生还。

\textbf{14.} 与飞行轨迹平行的斜坡可使 $v_\perp$
较小，所以双腿吸收的只是几米，
而不是整座动能之山。

\textbf{15.} $H'_0 = v_0'^2/2g = 576/19.62 \approx \qty{29.4}{m}$。

\textbf{16.} $mgH' - f\,\dfrac{H'}{\sin\qty{35}{\degree}} =
\tfrac12 m v_0'^2$。

\textbf{17.} $H' = \dfrac{\tfrac12 \times 70 \times 576}
{686.7 - 80/0.574} = \dfrac{\num{20160}}{547} \approx \qty{37}{m}$。

\textbf{18.} $36.8/29.4 \approx 1.25$：附加 $25\%$。

\textbf{19.} 对于 \qty{60}{kg}：$H' = \num{17280}/(588.6 - 139.5)
\approx \qty{38.5}{m}$。\emph{较轻}的跳跃者需要更高的
跳台：摩擦力的代价 $fL$ 不变，但她的能量预算
$mgH'$ 较小。

\textbf{20.} 宣布建造 \qty{37}{m} 高的跳台，以达到 \qty{24}{m/s} 的
起跳速率——并承认其中约四分之一，即约 \qty{7}{m} 的
混凝土，是摩擦力对这一无摩擦理想所附加的代价。
\end{solution}
